\usepackage[spanish, es-tabla]{babel} %Idioma

% Codificación
\usepackage[utf8]{inputenc}
\usepackage[T1]{fontenc}
\usepackage{qrcode}

% Paquetes matemáticos
\usepackage{amsmath, amssymb, physics, bigints, cancel}
\usepackage{stmaryrd}

% Colores y cuadros
\usepackage[most]{tcolorbox}
\usepackage[absolute,overlay]{textpos}
\setlength{\TPHorizModule}{1cm}
\setlength{\TPVertModule}{1cm}

\newtcolorbox{infobox}[2][blue]{
    colback=#1!5!white,
    colframe=#1!75!black,
    title={#2},
    fonttitle=\bfseries,
    coltitle=black,
    boxrule=0.8pt,
    arc=4pt,
    left=6pt,
    right=6pt,
    top=6pt,
    bottom=6pt,
    breakable
}

% TikZ para diagramas y circuitos
\usepackage{tikz}
\usetikzlibrary{patterns, decorations.pathreplacing}
\usepackage{circuitikz}

% Figuras
\usepackage{graphicx}
\usepackage{float} % no siempre necesario, pero puede ayudar

% Tablas pro
\usepackage{booktabs}

% Listados de código
\usepackage{listings}

% Bash personalizado
\lstdefinelanguage{custombash}{
    language=bash,
    morekeywords={conda, activate, mkdir, cd, cmake, make, git, sudo},
    sensitive=true
}

\lstdefinestyle{bashStyle}{
    language=custombash,
    basicstyle=\ttfamily\small,
    keywordstyle=\color{teal}\bfseries,
    commentstyle=\color{gray}\itshape,
    stringstyle=\color{orange},
    identifierstyle=\color{black},
    numbers=left,
    numberstyle=\tiny\color{gray},
    stepnumber=1,
    numbersep=10pt,
    backgroundcolor=\color{white},
    showspaces=false,
    showstringspaces=false,
    breaklines=true,
    frame=single,
    captionpos=b
}

% Python
\lstdefinestyle{pythonStyle}{
    language=Python,
    basicstyle=\ttfamily\small,
    keywordstyle=\color{blue}\bfseries,
    commentstyle=\color{gray}\itshape,
    stringstyle=\color{orange},
    identifierstyle=\color{black},
    numbers=left,
    numberstyle=\tiny\color{gray},
    stepnumber=1,
    numbersep=10pt,
    backgroundcolor=\color{white},
    showspaces=false,
    showstringspaces=false,
    breaklines=true,
    frame=single,
    captionpos=b
}

% C++
\lstdefinestyle{cppStyle}{
    language=C++,
    basicstyle=\ttfamily\small,
    keywordstyle=\color{blue}\bfseries,
    commentstyle=\color{gray}\itshape,
    stringstyle=\color{red},
    numbers=left,
    numberstyle=\tiny\color{gray},
    stepnumber=1,
    numbersep=10pt,
    backgroundcolor=\color{white},
    showspaces=false,
    showstringspaces=false,
    breaklines=true,
    frame=single,
    captionpos=b
}

% Fuente doble negrita
\usepackage{bbold}

% Íconos FontAwesome
\usepackage{fontawesome5}

% Para texto de ejemplo
\usepackage{lipsum}

% Comando para tachado en color
\newcommand{\colorcancel}[2]{\textcolor{#1}{\cancel{\textcolor{black}{#2}}}}

% Comando para círculos con texto
\newcommand*\circled[1]{\tikz[baseline=(char.base)]{
  \node[shape=circle,draw,inner sep=1.5pt] (char) {#1};}}

% Pie de página
\setbeamertemplate{footline}[frame number]

% Justificado
\usepackage{ragged2e} 
\usepackage{etoolbox}
\apptocmd{\itemize}{\justifying}{}{}
\apptocmd{\enumerate}{\justifying}{}{}
